% Esquisse d'un programme (cote 119, copy 1): archive cote 119, site page 3 (PDF page 4), typescript page 1.
% typescript. Draft compared: EsquisseFr.pdf (text layer):1-70.
% AI-transcribed by Claude Opus 5.5 (claude-opus-5-5, adaptive thinking, effort high), Message Batches API, batch pilot (msgbatch_0189Nhf1xB615YWnztubtLfc), 2026-11-12; not yet checked by eye.
\scanpage{119/3}
\mspage{1}

ESQUISSE D'UN PROGRAMME

par A. Grothendieck

1. Comme la conjoncture actuelle rend de plus en plus illu-%
soires pour moi les perspectives d'un enseignement de recherche à
l'Université, je me suis résolu à demander mon admission au CNRS, pour
pouvoir consacrer mon énergie à développer des travaux et perspectives
dont il devient clair qu'il ne se trouvera aucun élève (ni même,
semble-t-il, aucun congénère mathématicien) pour les développer à ma
place.

En guise de document ``Titres et Travaux'', on trouvera à la
suite de ce texte la reproduction intégrale d'une esquisse, par thèmes,
de ce que je considérais comme mes principales contributions mathéma-%
tiques au moment d'écrire ce rapport, en 1972. Il contient également
une liste d'articles publiés à cette date. J'ai cessé toute publica-%
tion d'articles scientifiques depuis 1970. Dans les lignes qui suivent,
je me propose de donner un aperçu au moins sur quelques thèmes princi-%
paux de mes réflexions mathématiques depuis lors. Ces réflexions se
sont matérialisées au cours des années en deux volumineux cartons de
notes manuscrites, difficilement déchiffrables sans doute à tout autre
qu'à moi-même, et qui, après des stades de décantations successives,
attendent leur heure peut-être pour une rédaction d'ensemble tout au
moins provisoire, à l'intention de la communauté mathématique. Le
terme ``rédaction'' ici est quelque peu impropre, alors qu'il s'agit
bien plus de développer des idées et visions multiples amorcées au
cours de ces douze dernières années, en les précisant et les appro-%
fondissant, avec tous les rebondissements imprévus qui constamment
accompagnent ce genre de travail - un travail de découverte donc, et
non de compilation de notes pieusement accumulées. Et je compte bien,
dans l'écriture des ``Réflexions Mathématiques'' commencée depuis fé-%
vrier 1983, laisser apparaître clairement au fil des pages la démarche
de la pensée qui sonde et qui découvre, en tâtonnant dans la pénombre
bien souvent, avec des trouées de lumière subites quand quelque tenace

\margin{3}
